% GENERATED FILE. Edit canonical lessons, then run npm run generate:books.
% canonical-source-hash: fnv1a64:8ec989ad3a6eee0e
% canonical-lessons: FR-C460-meteo, FR-C460-previsions, FR-C460-nuageux, FR-C460-ensoleille, FR-C460-thermometre

\chapter{Le bulletin météo}
\label{ch:fr-le-bulletin-meteo}
\hlchaptermodality{460}

\begin{chapteropening}
\textbf{By the end of this chapter:} \emph{I can follow a weather report: the weather forecast, the outlook, cloudy, sunny and the thermometer.}

Complete the last lesson of chapter 460: I can follow a weather report: the weather forecast, the outlook, cloudy, sunny and the thermometer.
\end{chapteropening}

\section[la météo]{la météo — the weather forecast; the weather}

\label{lesson:FR-C460-meteo}



\begin{quote}
Before the new one: say the French for the stag, then the French for the doe.
\end{quote}

\subsection*{You'll want to know: la météo}
\textbf{la météo} — \textquotedblleft{}the weather forecast; the weather\textquotedblright{}.

\begin{quote}
Selon la météo, il va faire beau demain.
\end{quote}

\begin{quote}
\emph{According to the forecast it is going to be nice tomorrow.}
\end{quote}

Short for \emph{la météorologie}.

\begin{cousinweb}
\emph{météo} is short for \emph{météorologie}, a learned borrowing from Greek \emph{meteōrología}, from \emph{metéōros} 'high in the air' (EtymDB lists the short form).
\end{cousinweb}

\begin{grammarlens}[title={What you've built}]
One more word for telling, explaining and describing at B1.
\end{grammarlens}

\subsection*{Your turn}
\begin{itemize}
\raggedright
  \item \emph{Say it:} \emph{la météo}
  \item \emph{Say it:} \emph{la météo}, once more
  \item \emph{Say it:} a sentence of your own with \emph{la météo}
  \item \emph{Recall:} say the French for the stag, then the French for the doe, then say \emph{la météo} again
\end{itemize}

\begin{tcolorbox}[breakable,colback=teal!4,colframe=teal!35!black,title={Before you move on}]
What does la météo mean? (\textquotedblleft{}The weather forecast; the weather\textquotedblright{}.) Say it once more.
\end{tcolorbox}
\section[les prévisions]{les prévisions — the forecast, the outlook}

\label{lesson:FR-C460-previsions}



\begin{quote}
Before the new one: say the French for the doe, then the French for the weather forecast.
\end{quote}

\subsection*{You'll want to know: les prévisions}
\textbf{les prévisions} — \textquotedblleft{}the forecast, the outlook\textquotedblright{}.

\begin{quote}
Les prévisions pour le week-end ne sont pas bonnes.
\end{quote}

\begin{quote}
\emph{The forecast for the weekend is not good.}
\end{quote}

Mostly plural: \emph{les prévisions météo}.

\begin{cousinweb}
\emph{prévision} is built on \emph{prévoir} 'to foresee', from Latin \emph{praevidēre}, built on \emph{vidēre} 'to see' (Lewis \& Short). EtymDB does not list the noun.
\end{cousinweb}

\begin{grammarlens}[title={What you've built}]
One more word for telling, explaining and describing at B1.
\end{grammarlens}

\subsection*{Your turn}
\begin{itemize}
\raggedright
  \item \emph{Say it:} \emph{les prévisions}
  \item \emph{Say it:} \emph{les prévisions}, once more
  \item \emph{Say it:} a sentence of your own with \emph{les prévisions}
  \item \emph{Recall:} say the French for the doe, then the French for the weather forecast, then say \emph{les prévisions} again
\end{itemize}

\begin{tcolorbox}[breakable,colback=teal!4,colframe=teal!35!black,title={Before you move on}]
What does les prévisions mean? (\textquotedblleft{}The forecast, the outlook\textquotedblright{}.) Say it once more.
\end{tcolorbox}
\section[nuageux]{nuageux — cloudy}

\label{lesson:FR-C460-nuageux}



\begin{quote}
Before the new one: say the French for the weather forecast, then the French for the forecast.
\end{quote}

\subsection*{You'll want to know: nuageux}
\textbf{nuageux} — \textquotedblleft{}cloudy\textquotedblright{}.

\begin{quote}
Le ciel est nuageux, mais il ne pleut pas.
\end{quote}

\begin{quote}
\emph{The sky is cloudy, but it is not raining.}
\end{quote}

Feminine \emph{nuageuse}.

\begin{cousinweb}
\emph{nuageux} is built on \emph{nuage} 'cloud' (EtymDB lists it), from the old word \emph{nue} 'cloud', which goes back to Latin \emph{nūbēs}.
\end{cousinweb}

\begin{grammarlens}[title={What you've built}]
One more word for telling, explaining and describing at B1.
\end{grammarlens}

\subsection*{Your turn}
\begin{itemize}
\raggedright
  \item \emph{Say it:} \emph{nuageux}
  \item \emph{Say it:} \emph{nuageux}, once more
  \item \emph{Say it:} a sentence of your own with \emph{nuageux}
  \item \emph{Recall:} say the French for the weather forecast, then the French for the forecast, then say \emph{nuageux} again
\end{itemize}

\begin{tcolorbox}[breakable,colback=teal!4,colframe=teal!35!black,title={Before you move on}]
What does nuageux mean? (\textquotedblleft{}Cloudy\textquotedblright{}.) Say it once more.
\end{tcolorbox}
\section[ensoleillé]{ensoleillé — sunny}

\label{lesson:FR-C460-ensoleille}



\begin{quote}
Before the new one: say the French for the forecast, then the French for cloudy.
\end{quote}

\subsection*{You'll want to know: ensoleillé}
\textbf{ensoleillé} — \textquotedblleft{}sunny\textquotedblright{}.

\begin{quote}
La terrasse est ensoleillée l'après-midi.
\end{quote}

\begin{quote}
\emph{The terrace is sunny in the afternoon.}
\end{quote}

\emph{Une journée ensoleillée} — a sunny day.

\begin{cousinweb}
\emph{ensoleillé} is built on \emph{soleil} 'sun', which goes back to Vulgar Latin \emph{sōliculus}, a small \emph{sōl} 'sun' (Lewis \& Short). EtymDB does not list it.
\end{cousinweb}

\begin{grammarlens}[title={What you've built}]
One more word for telling, explaining and describing at B1.
\end{grammarlens}

\subsection*{Your turn}
\begin{itemize}
\raggedright
  \item \emph{Say it:} \emph{ensoleillé}
  \item \emph{Say it:} \emph{ensoleillé}, once more
  \item \emph{Say it:} a sentence of your own with \emph{ensoleillé}
  \item \emph{Recall:} say the French for the forecast, then the French for cloudy, then say \emph{ensoleillé} again
\end{itemize}

\begin{tcolorbox}[breakable,colback=teal!4,colframe=teal!35!black,title={Before you move on}]
What does ensoleillé mean? (\textquotedblleft{}Sunny\textquotedblright{}.) Say it once more.
\end{tcolorbox}
\section[le thermomètre]{le thermomètre — the thermometer}

\label{lesson:FR-C460-thermometre}



\begin{quote}
Before the new one: say the French for cloudy, then the French for sunny.
\end{quote}

\subsection*{You'll want to know: le thermomètre}
\textbf{le thermomètre} — \textquotedblleft{}the thermometer\textquotedblright{}.

\begin{quote}
Ce matin, le thermomètre indique moins cinq.
\end{quote}

\begin{quote}
\emph{This morning the thermometer reads minus five.}
\end{quote}

Also for a fever.

\begin{cousinweb}
\emph{thermomètre} is a learned formation from Greek \emph{thermós} 'hot' and \emph{métron} 'measure' (EtymDB lists the word).
\end{cousinweb}

\begin{grammarlens}[title={What you've built}]
One more word for telling, explaining and describing at B1.
\end{grammarlens}

\subsection*{Your turn}
\begin{itemize}
\raggedright
  \item \emph{Say it:} \emph{le thermomètre}
  \item \emph{Say it:} \emph{le thermomètre}, once more
  \item \emph{Say it:} a sentence of your own with \emph{le thermomètre}
  \item \emph{Recall:} say the French for cloudy, then the French for sunny, then say \emph{le thermomètre} again
\end{itemize}

\begin{tcolorbox}[breakable,colback=teal!4,colframe=teal!35!black,title={Before you move on}]
What does le thermomètre mean? (\textquotedblleft{}The thermometer\textquotedblright{}.) Say it once more.
\end{tcolorbox}
